% page 405 du fac-simile (image p-404 du scan)
% bloc 165.1 x 233.3 mm ; marge gauche 36.9 mm
\pgc{15.240mm}{0.000mm}{
\l{\cel{54}397}
\l{}
\l{\sou{nostalgio}. Langoro quan produktas la regretado di la nasko-lando.}
\l{\cel{3}EFIRS.}
\l{}
\l{\sou{nostoko.}\cel{2}(\sou{historio naturala}). Algo filamentoza, envelopita da}
\l{\cel{3}gaino gelatinatra. - L.\cel{1}\sou{nostoc}. - DEFS.}
\l{}
\l{\sou{notar}.\cel{2}(\sou{trans.}) Skribar to di quo onu volas konservar la men-}
\l{\cel{3}ciono, la indiko. - DEFIRS.}
\l{}
\l{\sou{notario}. Oficiero publika, establisita por redaktar e recevar ol-}
\l{\cel{3}ti ek la akti, kontrati, e c., ad olqui la privati volas do-}
\l{\cel{3}nar karaktero di autentikeso. DEFIRS.}
\l{}
\l{\sou{notico}. Skriburo destinata ad informar pri punto partikulara di}
\l{\cel{3}historio, di cienco, e c.\cel{2}- DEFIS.}
\l{}
\l{\sou{notifikar}. (\sou{trans.})\cel{2}Igar ulu savar, uzante por lo la formo}
\l{\cel{3}oficala. - DEFIS.}
\l{}
\l{\sou{notora}.\cel{2}Qua konocesas, savesas, da publiko. DEFIRS.}
\l{}
\l{\sou{nova}. I. Qua aparas unesmafoye, od aparis erste de poka tempo.}
\l{\cel{3}- II. Qua aparas pos altra, quan lu remplasas. - III.}
\l{\cel{3}( nuva): Quan onu ne jakomencis konsumar, o quan on jus}
\l{\cel{3}komencis konsumar. - DEFIRS.}
\l{}
\l{\sou{novelo.}\cel{1}(\sou{literaturo)}.\cel{2}Naraco (kelka-pagina) pri aventuro intere-}
\l{\cel{3}siva. - DEFIRS.}
\l{}
\l{\sou{novembro}. La monato dek-ed-unesma di la yaro. - DEFIRS.}
\l{}
\l{\sou{novico}. La persono,qua,jus diveninte reguliero (che la katoliki),}
\l{\cel{3}subisas probado dum ula tempo ante profesar. - (\sou{anke metaf.})}
\l{\cel{3}DEFIS.}
\l{}
\l{nu !\cel{5}Interjeciono qua : (l) segun la pronunco-tono, expresas}
\l{\cel{3}astoneso o konsento ; (2) preiras afirmajo quan onu qualifikas}
\l{\cel{3}kom stranja, neexpektita. - DR.}
\l{}
\l{\sou{nuanco}. I. Singla di la toni di un koloro, de la maxim klara}
\l{\cel{3}til la maxim obskura. - II. (\sou{metaf}.) Difero di grado, preske}
\l{\cel{3}nesentebla, neperceptebla, qua distingas kozo de olti qui esas}
\l{\cel{3}lua-vicina. - DEFR.}
\l{}
\l{\sou{nubo}.\cel{2}(\sou{meteor.})\cel{2}Amaso de aquo-vaporo, qua trublas la diafaneso}
\l{\cel{3}di atmosfero, e finas per divenar pluvo-guti.- EIS.}
\l{}
\l{\sou{nuco}.\cel{2}(\sou{bot.)}\cel{2}Frukto di arboro ek la familio \textquotedbl{}juglandei\textquotedbl{}, di}
\l{\cel{3}qua la ligno uzesas da la ebenisti, moblifisti. - L.\cel{1}\sou{juglans}.}
\l{}
\l{nucelo. (\sou{biol.}) Parto di la ovulo, qua envelopas la sako embrio-}
\l{\cel{3}nala e qua tegesas da la tegumenti di la ovulo. -}
}
